% ============================================================
% 论文 D Outline — 约束增强的解耦规划框架
% 目标期刊: Robotics and Autonomous Systems (Q2, IF 5.2)
% 语言: 英文（此文件为中文行文逻辑大纲，正式论文用英文撰写）
% ============================================================

\documentclass[a4paper,12pt]{article}
\usepackage[UTF8]{ctex}
\usepackage{geometry,enumitem,amsmath,booktabs}
\geometry{margin=2.5cm}
\setlist{nosep}

\title{论文D 行文逻辑大纲\\[4pt]
\large Constraint-Enriched Decoupled Planning:\\
Bridging the Path-Speed Information Gap\\for Multi-Obstacle Autonomous Driving}
\author{廖嘉旺}
\date{2026年5月}

\begin{document}
\maketitle

% ============================================================
\section*{核心叙事（一句话）}
% ============================================================

解耦规划的 path$\to$speed 信息断裂是已知问题；
现有方案要么推倒重来（联合规划，计算翻倍），要么改目标函数（soft cost，信息隐式）；
我们提出第三条路——\textbf{显式约束传递通道}，计算代价几乎不增加，
在生产级 Apollo 11.0 上实车验证有效。

贡献声明（3 条，但本质是 1 个框架）：
\begin{enumerate}
  \item 提出约束增强解耦规划框架（Constraint-Enriched Decoupled Planning, CEDP），
        通过显式信息通道恢复时空协调能力；
  \item 证明解耦规划的不可行性条件（Proposition）及通道恢复可行性的保证；
  \item 在 Apollo 11.0 生产级代码上实现并通过实车验证。
\end{enumerate}

% ============================================================
\section{Introduction}
% ============================================================

行文逻辑链：

\begin{enumerate}
  \item \textbf{解耦是工业主流。}
        Apollo EM Planner、JD ROVER 5.0、Autoware 均采用 path-speed 解耦。
        原因：计算效率高（3D$\to$2$\times$2D）、模块化好、工程可维护。
  \item \textbf{但解耦有结构性盲区。}
        path 阶段不知道时间$\to$动态障碍物只能做静态投影$\to$过度保守或漏判。
        speed 阶段拿到固定 path$\to$只能刹车，不能绕。
        用一个具体场景图（Figure 1）3 秒说清问题。
  \item \textbf{现有方案的局限。}
        \begin{itemize}
          \item 联合规划（IA Planner, 胡杰 2023/2025）：推倒重来，计算 76--117\,ms vs 解耦 35\,ms。
          \item 博弈选择（Zhang 2023 DP-RCG）：并行枚举 $O(mn)$，架构改动大。
          \item 框架内 soft cost（Han 2026）：信息隐式传递，无可行性保证。
        \end{itemize}
  \item \textbf{我们的 insight。}
        不需要替换架构，只需要在两个阶段之间建一条\textbf{显式约束通道}。
        通道回答三个问题：传什么（which obstacles）、空间怎么传（where）、时间怎么传（when）。
  \item \textbf{贡献声明。}（见上）
\end{enumerate}

% ============================================================
\section{Problem Formulation \& Infeasibility Analysis}
% ============================================================

\subsection{解耦架构形式化}

\begin{itemize}
  \item Stage 1: Path QP — 输入：静态/准静态障碍物边界 $[l_{\min}(s), l_{\max}(s)]$；
        输出：最优路径 $l^*(s)$。
  \item Stage 2: Speed QP — 输入：固定路径 $l^*(s)$ + ST 边界；
        输出：速度曲线 $s^*(t)$。
  \item 接口：Stage 1 $\to$ Stage 2 传递的是 $l^*(s)$，\textbf{不含时间信息}。
\end{itemize}

\subsection{信息损失刻画}

定义联合可行域 $\mathcal{F}_{\text{joint}}$ vs 解耦可行域 $\mathcal{F}_{\text{dec}}$。
当动态障碍物存在时，$\mathcal{F}_{\text{dec}} \subseteq \mathcal{F}_{\text{joint}}$，
且存在场景使 $\mathcal{F}_{\text{dec}} = \emptyset$ 而 $\mathcal{F}_{\text{joint}} \neq \emptyset$。

\subsection{Proposition 1: Decoupling Infeasibility}

\textbf{命题：} 给定障碍物集 $\mathcal{O}$，若存在动态障碍物 $o_i$ 使得
其预测占用区域在静态投影下覆盖了所有横向可行带，即
$\forall s \in [s_a, s_b],\; \bigcup_{t} \text{Proj}_{SL}(o_i, t) \supseteq [l_{\min}(s), l_{\max}(s)]$，
则 $\mathcal{F}_{\text{dec}} = \emptyset$。

\textbf{证明思路：} Path 阶段无 $\tau(s)$ 信息 $\to$ 必须对所有时刻取并集作为占用
$\to$ 占用膨胀 $\to$ 可行带被压缩为空。

\textbf{对比：} 约束增强规划器知道 $\tau(s)$，只查询 $o_i$ 在 $t=\tau(s)$ 时刻的位置
$\to$ 占用不膨胀 $\to$ 可行带保留。

% ============================================================
\section{Constraint-Enriched Decoupled Planning}
% ============================================================

\textbf{按信息流组织，不按模块组织。}

\subsection{框架概览（Figure 2: 架构图）}

一张图展示：
\[
\text{Perception} \to \underbrace{\text{Threat Filter} \to \text{Spatial Decision} \to \text{Path QP}}_{\text{Stage 1 (enriched)}}
\xrightarrow{\mathcal{T} = \{(s_k, \tau_k, \delta_k)\}}
\underbrace{\text{Corridor Construction} \to \text{Speed QP}}_{\text{Stage 2 (enriched)}}
\]

通道 $\mathcal{T}$ 携带：纵向位置 $s_k$、到达时间 $\tau_k$、安全裕度 $\delta_k$。

\subsection{What to Pass: 威胁度过滤}

不是所有障碍物都需要进入通道。定义多因子威胁度：
\[
\Theta_i = \sum_{j=1}^{5} w_j f_j(o_i)
\]
五因子：TTC、横向重叠、相对速度、障碍物类型、交互密度。

只有 $\Theta_i > \Theta_{\text{thr}}$ 的障碍物进入约束通道。
这是过滤器，不是贡献本身——降低通道负载。

\subsection{Where to Pass: 最大间隙空间决策}

对进入通道的障碍物组，确定最优通行带：

\textbf{Theorem (Max-Gap Optimality):}
对 $n$ 个纵向重叠障碍物，扫描线分组后，最优通行带宽等于
排序边界中相邻间隙的最大值，可在 $O(n \log n)$ 时间内求得。

\textbf{证明：} 引理（连续划分性）$\to$ 定理（最大间隙 = 最优宽度）。

输出：通行带边界 $[l_{\min}^*(s), l_{\max}^*(s)]$ 送入 Path QP。

\subsection{When to Pass: ST 走廊时间约束}

基于 Path QP 输出的 $l^*(s)$ 和运动学模型，计算到达时间函数 $\tau(s)$。

构造 ST 走廊：
\[
\Omega = \{(s,t) \mid g_{p^*}(s) \geq W_{\text{ego}},\; t \geq \tau(s)\}
\]

将 $\Omega$ 作为约束注入 Speed QP。

\textbf{Proposition 2 (Feasibility Recovery):}
若 $\Omega \neq \emptyset$ 且为凸集，则 Speed QP 在 $\Omega$ 内必有可行解。

（注：诚实标注为 Proposition，不装 Theorem。前提条件"凸"需要无禁行切片。）

\subsection{Verification Loop（可选反馈）}

Speed QP 求解后，检查 $s^*(t)$ 是否满足 $t \geq \tau(s)$。
若违反，用更新的速度曲线重新计算 $\tau'(s)$，裁剪走廊，再跑一次 Speed QP。
最多迭代一次，计算代价 $\leq 2\times$ 单次 QP。

% ============================================================
\section{Theoretical Analysis}
% ============================================================

汇总理论结果（可合并到 Section 2\&3，也可独立成节，视篇幅定）：

\begin{table}[h]
\centering
\begin{tabular}{lll}
\toprule
结果 & 类型 & 内容 \\
\midrule
Proposition 1 & 不可行性 & 解耦在动态覆盖条件下必然失效 \\
Theorem & 最优性 & 最大间隙 $O(n\log n)$ \\
Proposition 2 & 可行性恢复 & $\Omega \neq \emptyset \Rightarrow$ Speed QP 有解 \\
\bottomrule
\end{tabular}
\end{table}

复杂度对比表：
\begin{table}[h]
\centering
\begin{tabular}{lcc}
\toprule
方法 & 时间复杂度 & 架构改动 \\
\midrule
原版解耦 (Apollo) & $O(n)$ per stage & 无 \\
联合规划 (IA/NMPC) & $O(n^2)$--$O(n^3)$ & 全部替换 \\
博弈选择 (Zhang 2023) & $O(mn)$ 枚举 & 并行生成模块 \\
\textbf{CEDP (ours)} & $O(n \log n)$ + 2$\times$QP & 仅加通道 \\
\bottomrule
\end{tabular}
\end{table}

% ============================================================
\section{Implementation on Apollo 11.0}
% ============================================================

\begin{itemize}
  \item 修改模块：\texttt{PathBoundsDecider}（注入威胁度+间隙决策）、
        \texttt{SpeedBoundsDecider}（接收走廊约束）、新增 \texttt{ConstraintChannel}。
  \item 代码侵入度：新增 $\sim$800 行 C++，修改 $\sim$200 行，不改 QP 求解器本身。
  \item 与原架构兼容：通道为空时退化为原版 Apollo 行为（graceful degradation）。
\end{itemize}

% ============================================================
\section{Experiments}
% ============================================================

\subsection{仿真实验}

\begin{itemize}
  \item 平台：Apollo 11.0 Dreamview，10\,Hz 规划周期。
  \item 场景：15--20 个配置（5 基础场景 $\times$ 3--4 参数变体：障碍物速度、数量、间距）。
  \item 每个配置跑 10 次，报告 mean $\pm$ std。
  \item Metrics：通过率、平均巡航速度、最小障碍物距离、纵向 jerk、全流水线耗时。
  \item Baseline：原版 Apollo EM Planner。
  \item 额外对比（如可行）：简化版 MPC 或 Zhang 2023 逻辑复现。
\end{itemize}

\subsection{实车验证}

\begin{itemize}
  \item 平台：Apollo 实车，封闭场地。
  \item 场景：2--3 个（静态锥桶绕行、动态行人/小车穿越）。
  \item 交付物：规划曲线（曲率/速度/加速度 vs 时间）+ 实际执行曲线 + 场地照片/截图。
  \item 对比：同场景下原版 Apollo 的表现（停车/绕不过去）。
\end{itemize}

\subsection{消融实验}

关闭通道各维度：
\begin{itemize}
  \item 无威胁度过滤（所有障碍物进通道）$\to$ 计算开销增加
  \item 无间隙决策（用 Apollo 原版 nudge）$\to$ 通过率下降
  \item 无 ST 走廊（不传时间约束）$\to$ 动态场景失效
  \item 无 verification loop $\to$ 边界 case 退化
\end{itemize}

\subsection{计算开销分析}

全流水线 breakdown：威胁度计算 / 间隙决策 / Path QP / 走廊构造 / Speed QP。
证明总耗时仍在 100\,ms 规划周期内。

\subsection{失败边界}

参数扫描：在什么障碍物密度/速度下走廊变空（$\Omega = \emptyset$）？
此时系统 graceful degradation 为原版 Apollo 行为（停车）。

% ============================================================
\section{Related Work}
% ============================================================

放在实验之后（RAS 接受这种结构）。分三类：

\begin{enumerate}
  \item \textbf{联合规划（替换架构）：}
        IA Planner (Yang 2021), 胡杰 (2023, 2025), Tong (2025), 乔应 (2025)。
        共同点：推倒解耦，计算代价翻倍。
  \item \textbf{框架内修补（保留架构）：}
        Han (2026, soft cost), Zhang (2023, 博弈选择), Tong (2022, 轨迹修复)。
        我们的定位：显式硬约束通道 vs 隐式 soft cost / 枚举选择。
  \item \textbf{解耦架构的工业实践：}
        Apollo EM Planner (Fan 2018), JD ROVER (Wang 2022), Meng (2019)。
        说明解耦是主流，我们的工作对工业有直接价值。
\end{enumerate}

% ============================================================
\section{Conclusion}
% ============================================================

\begin{itemize}
  \item 总结：提出 CEDP 框架，通过显式约束通道恢复解耦规划的时空协调能力。
  \item 关键结果：通过率 0\%$\to$100\%，计算代价仅增加 $<$15\%，实车验证可行。
  \item 局限：当前验证限于 Apollo 平台，走廊凸性依赖无禁行切片假设。
  \item 未来：扩展到更多平台（Autoware）、引入学习型威胁度模型、大规模 benchmark 评测。
\end{itemize}

% ============================================================
\section*{时间线}
% ============================================================

\begin{tabular}{ll}
\toprule
时间 & 任务 \\
\midrule
5月下旬 & Outline 定稿，开始写 Section 1--3 \\
6月 & 补仿真场景（15--20个），写 Section 4--5 \\
7月 & 实车验证（2--3天），写 Section 6--7，全文 polish \\
8月初 & 投稿 Robotics and Autonomous Systems \\
\bottomrule
\end{tabular}

\end{document}
